% ch03 -- Träger und Information
\label{ch:traeger}
% Quelldokumente: A272

\epigraph{Die populäre Kurzfassung übersteigt, was Landauers Beweis zeigt.}{A272}

\section{Die Fragestellung}

Seit Landauers Arbeit von 1961 wird behauptet: Jedes gelöschte Bit muss mindestens
$k_BT\ln 2$ Wärme erzeugen. Diese Verallgemeinerung geht über das hinaus, was
Landauers Beweis zeigt. Dok.\,A272 führt die Textkritik; die physikalische Substanz
steht in A271.

\section{Zwei Definitionen}

\textbf{Information} ist eine abstrakte Größe. Ein Bit ist eine Einheit der
Information. Information hat keine Masse, keine Energie, keine Entropie und keine
physikalische Lokalisation.

\textbf{Träger} ist ein physikalisches System, das Information repräsentiert.
Der Träger hat Energie, Entropie und weitere physikalische Eigenschaften.

Diese Trennung ist eine Setzung. Sie wird nicht bewiesen; sie legt fest, worüber
geredet wird. Der gesamte Ertrag hängt daran, dass sie durchgehalten wird.

\section{Wörterbuch A271/A272}

Beide Dokumente treffen dieselbe Unterscheidung in verschiedenen Wörtern:

\begin{center}
\begingroup\small
\begin{tabular}{@{}p{2.8cm}p{2.8cm}p{2.0cm}@{}}
\toprule
\textbf{A271 (Buchhaltung)} & \textbf{A272 (Textkritik)} & \textbf{Ebene} \\
\midrule
Gebiet im Phasenraum & Träger & physikalisch \\
Beschreibung & Information & abstrakt \\
Gebietsreduktion $\{A,B\}\to Z$ & RESTORE TO ONE & physikalisch \\
Nicht-Bijektivität & Zustandsraumverkleinerung & physikalisch \\
Beschreibungsentscheidung & interpretative Löschung & abstrakt \\
\bottomrule
\end{tabular}
\endgroup
\end{center}

\section{Energie hängt am Träger}

\begin{theorem}
Einem Informationsbit kann keine Energie zugeordnet werden, nur seinem Träger.
\end{theorem}

\textit{Beweis.} Angenommen, ein Bit hätte intrinsische Energie. Diese müsste
unabhängig von der physikalischen Realisierung sein. Ein Bit $= 1$ kann jedoch
durch CMOS (5\,V), durch einen magnetischen Spin oder ein Photon realisiert werden
--- bei identischem Informationsgehalt, aber verschiedenen Energien. Also ist
Energie keine Eigenschaft der Information, sondern des Trägers.

Das Argument ist ein Mehrfachrealisierungs-Argument und trägt genau so weit,
wie die Prämisse gilt: Ein Bit bleibt dasselbe Bit bei Wechsel des Trägers.
Wo diese Invarianz bricht (Kap.\,\ref{ch:offene_fragen}), ist der Beweis
zu ergänzen.

\section{Hypostasierung als sprachlicher Mechanismus}

Wenn man von \enquote{dem Bit} spricht, behandelt man die abstrakte Kategorie
so, als hätte sie eine eigene physikalische Existenz --- das ist
\textbf{Hypostasierung}. Die Digitaltechnik begünstigt diesen Fehler strukturell:
Das Wort \enquote{bit} bezeichnet gleichzeitig die Hardwareeinheit (Träger) und
die Informationseinheit (abstrakt). Ein Sprache ohne diese Doppeldeutigkeit würde
den Fehler schwerer machen.

\section{Reichweite und Grenzen der Aussage}

Was hier nicht bestritten wird: Landauers Rechnung ist korrekt für die
Operation, die er betrachtet --- die physikalische RESTORE-TO-ONE-Operation auf
einem Träger im thermischen Gleichgewicht. Bestritten wird die Zuständigkeit
dieses Ergebnisses für abstrakte Information.
